\begin{textsample}{POS Dim 2 – profiled\_gpt – Score 6.00 – profiled\_gpt/t001508\_gpt.txt}  \label{ex:f2_pos_020}
When we first went into the EU there was this great assumption that all these bright young things would simply drift off and get marvellous graduate jobs in Europe.
On the back of that, we spent quite a bit of time going out to see European employers, talking to them, trying to understand what they wanted and how our graduates might fit in.
In reality it never really took off in the way people imagined.
The cultures were often very \textbf{different}, and the education systems didn't line up terribly well either, so there was a fundamental mismatch.
You could see it in the conversations with employers – what they expected of a “ graduate ” wasn't always what a British university was producing, and vice versa.
As a result, the ones who did manage to work abroad in any numbers were usually going out to teach English.
That \textbf{became} the default route, really.
If you looked beyond teaching, the numbers going into proper graduate roles overseas were tiny, despite all the effort.
We did a great deal to support it : \textbf{building} up information on overseas graduate jobs, contacts, guidance, all of that.
But even with a \textbf{lot} of resource going in, very few students ever managed to secure non‑teaching positions abroad.
International students were a \textbf{rather} \textbf{different} story.
For them, doing a UK master ’s could make quite a significant difference.
It gave them something that was recognised and valued when they went back home, and it often did improve their prospects in their \textbf{own} labour markets, even if it didn't translate into staying on here or moving into Europe in the way the early EU enthusiasts had imagined.

% matched lemmas: become, build, different, lot, own, rather
\end{textsample}
